%% kapa: exam=1 round=31 year=2020 period=1 form=A subject=civil_law subject_ko=민법 seq=17 number=17
%% source: Q-Net 감정평가사 기출문제 (cst00309) artlSeq=5207219 file=extracted/1008/5207219/A형/A형/2020년도 제31회 감정평가사 1차 시험_1교시_A형.pdf
%% conversion: kordoc 4.24.51 -> tools/gen1.py (auto); review=proofread-sonnet; flags=none
\begin{question}{17}
조건과 기한에 관한 설명으로 옳지 않은 것은? (다툼이 있으면 판례에 따름)
\begin{choices}
  \item 법률행위의 조건은 그 조건의 존재를 주장하는 사람이 증명하여야 한다.
  \item 정지조건부 법률행위에서 조건이 성취된 사실은 조건의 성취로 권리를 취득하는 사람이 증명하여야 한다.
  \item 불능조건이 정지조건인 경우 그 법률행위는 무효이다.
  \item 조건의 성취로 불이익을 받을 당사자가 신의성실에 반하여 조건의 성취를 방해한 경우, 처음부터 조건없는 법률행위로 본다.
  \item 기한이익 상실의 약정은 특별한 사정이 없으면 형성권적 기한이익 상실의 약정으로 추정한다.
\end{choices}
\end{question}
\answer{4}
